\section{Prime-residue pressure and the cost of an atypical profile}
\label{sec:residuepressure}
The weighted theorem has an arithmetic consequence beyond the
unweighted total count.
Keep a fixed even modulus \(q\), a nonempty reduced residue set
\(R\) with \(g_R=2\), and write \(k=|R|\).
For each unoriented Hamilton cycle \(C\) of \(G_N(q,R)\), put
\[
 n_r(C)=|\{ij\in E(C):i+j\equiv r\pmod q\}|,\qquad r\in R.
\]
These are edge counts, including repeated prime values if they
occur on different edges. Every cycle has exactly \(N\) edges,
so \(\sum_{r\in R}n_r(C)=N\).
For fixed weights \(w_r>0\), define
\[
 W_N(w)=\sum_{C\ {\rm Hamilton}}\prod_{r\in R}w_r^{\,n_r(C)},
 \qquad
 M(w)=\frac{\sum_{r\in R}w_r}{\varphi(q)}.
\]
The weights, modulus and allowed set are fixed before \(N\)
tends to infinity. No uniformity for growing moduli or
\(N\)-dependent weights is asserted.

\begin{theorem}[The full positive residue-pressure formula]
\label{thm:residuepressure}
Along even \(N\),
\[
 \log W_N(w)
 =N\left(\log\frac{N}{e\varphi(q)\log N}
                      +\log\sum_{r\in R}w_r\right)+o_w(N).
\]
In particular, if \(C_N\) is a uniform unoriented Hamilton cycle
of \(G_N(q,R)\) and \(X_N=(n_r(C_N)/N)_{r\in R}\), then for every
fixed \(\theta\in\mathbb R^R\),
\[
 \lim_{N\to\infty,\ 2\mid N}\frac1N
      \log\mathbb E\exp\bigl(N\langle\theta,X_N\rangle\bigr)
       =\log\left(\frac1k\sum_{r\in R}e^{\theta_r}\right).
\]
\end{theorem}

\begin{proof}
The positive-weight version of the signed Fourier model is
obtained by replacing the uniform average \(\rho_R\) by
\[
 \rho_w(b)=
       \frac{\sum_{r\in R}w_r\e(br/q)}{\sum_{r\in R}w_r}.
\]
For completeness, each matrix entry is a finite linear
combination of the same fixed-class prime convolution
operators used in Lemma~\ref{lem:filtermodel}.
The Chinese remainder computation there consequently gives
the exact reflection coefficient
\[
 w(\gamma)=\rho_w(b)\frac{\mu(c)}{\varphi(c)},\qquad
 \gamma=b/q+a/c,\quad (c,q)=1,
\]
with the same squarefree and primitive-conductor conditions
as in (38), and zero in every other case.
All phases are retained before taking eigenvalues.
A positive weighted average of unit complex numbers has
modulus one precisely when all the phases agree.
Its set of unit-modulus characters is therefore unchanged:
\(b=0,q/2\), because \(g_R=2\).
The nontrivial gap is the positive fixed number
\[
 \delta_w=
 1-\max\left\{\frac1{p_q-1},
       \max_{\substack{b\bmod q\\b\ne0,\ q/2}}|\rho_w(b)|
          \right\}>0.
\]
The finite-mask operator approximation in (39) is linear
and has only a fixed weight-dependent constant.
It therefore proves the corresponding spectral approximation,
not merely a pointwise residue heuristic.

Consider the true-prime weighted kernel
\[
 Q_w(i,j)=
 \frac{w_{\,i+j\bmod q}\log(i+j)}{M(w)N}
      \one_{\substack{i\ne j,\ i+j\ {\rm prime}\\
                       i+j\bmod q\in R}}.
\]
The fixed-progression prime number theorem gives its row
defect \(O_{q,w,B}((\log N)^{-B})\) for every fixed \(B\).
Proper prime powers and the diagonal are removed with the
same row-norm bound as in the filtered model.
Thus its normalized spectrum has only the constant and
parity exceptions and gap \(\delta_w-o(1)\).
Lemma~\ref{lem:scaling}, with any fixed \(B>4\), produces
\[
 P_w=\diag(u)Q_w\diag(u),\qquad
 P_w\one=\one,\quad
 \|u-\one\|_\infty=o(1).
\]
This kernel has a fixed nontrivial gap, the original true-prime
support and maximum entry \(O_{q,w}(\log N/N)\).
It satisfies Theorem~\ref{thm:weightedentropy}, giving
\[
 \log Z_{\rm H}(P_w)=-N+o_w(N).
\]
Every Hamilton-cycle weight changes by the common factor
\(\prod_vu_v^2=\exp(o_w(N))\), so the same logarithmic formula
holds for \(Z_{\rm H}(Q_w)\).

One must still remove the nonconstant logarithmic prime weight.
This cannot be justified by a uniform entrywise \(1+o(1)\)
comparison at the smallest sums.
Put \(D_w=M(w)N/\log N\). Exactly,
\[
 \prod_{ij\in E(C)}Q_w(i,j)
 =D_w^{-N}\left(\prod_{r\in R}w_r^{n_r(C)}\right)
             \prod_{ij\in E(C)}\frac{\log(i+j)}{\log N}.
\]
The logarithm of the last product is \(o(N)\), uniformly
over all such cycles.
Indeed, for sums at least \(N/(\log N)^4\), each logarithm
has modulus \(O(\log\log N/\log N)\).
Every smaller sum has both endpoints in the first
\(N/(\log N)^4\) vertices; a cycle has at most twice that
many edges incident with this set.
For each of these edges the logarithm has modulus
\(O(\log\log N)\), since the smallest permitted sum is at least 3.
The complete exceptional contribution is
\(O(N\log\log N/(\log N)^4)=o(N)\).
Consequently
\[
 Z_{\rm H}(Q_w)=D_w^{-N}W_N(w)\exp(o_w(N)).
\]
This proves the partition-function formula.
Divide by its specialization \(w_r=1\), and then put
\(w_r=e^{\theta_r}\), to obtain the moment-generating formula.
\end{proof}

Let \(\Delta_R=\{x\in[0,1]^R:\sum_rx_r=1\}\).
Use \(0\log0=0\) in
\[
 I_R(x)=\sum_{r\in R}x_r\log(kx_r),\qquad x\in\Delta_R.
\]
Outside this simplex set \(I_R=+\infty\).

\begin{theorem}[A full empirical residue large-deviation law]
\label{thm:residueldp}
The random vectors \(X_N\) satisfy the large-deviation principle
with speed \(N\) and rate \(I_R\).
Precisely, for every relatively open \(O\subseteq\Delta_R\)
and relatively closed \(F\subseteq\Delta_R\),
\[
 \begin{split}
 \liminf_{N\to\infty,\ 2\mid N}
       N^{-1}\log\mathbb P(X_N\in O)&\geq-\inf_{x\in O}I_R(x),\\
 \limsup_{N\to\infty,\ 2\mid N}
       N^{-1}\log\mathbb P(X_N\in F)&\leq-\inf_{x\in F}I_R(x).
 \end{split}
\]
Here \(\log0=-\infty\).
In particular, for each \(x\in\Delta_R\), if
\(B_\epsilon(x)=\{y\in\Delta_R:\|y-x\|_\infty<\epsilon\}\),
then both
\[
 \lim_{\epsilon\downarrow0}\liminf_N
       N^{-1}\log\mathbb P(X_N\in B_\epsilon(x)),
 \qquad
 \lim_{\epsilon\downarrow0}\limsup_N
       N^{-1}\log\mathbb P(X_N\in B_\epsilon(x))
\]
equal \(-I_R(x)\), with the inner limits along even \(N\).
\end{theorem}

\begin{proof}
We give the exponential-tilt proof to expose the quantifiers.
Write
\(\psi(\theta)=\log(k^{-1}\sum_re^{\theta_r})\).
Markov's inequality and Theorem~\ref{thm:residuepressure}
give, for each fixed \(\theta,x,\epsilon\),
\[
 \limsup_N N^{-1}\log\mathbb P(X_N\in B_\epsilon(x))
 \leq \psi(\theta)-\langle\theta,x\rangle
                                      +\epsilon\|\theta\|_1.
\]
If all \(x_r>0\), take \(\theta_r=\log(kx_r)\), for which
\(\psi(\theta)=0\) and \(\langle\theta,x\rangle=I_R(x)\).
If some coordinates vanish, set their \(\theta_r=-M\)
and use the same positive-coordinate values, then let
\(M\to\infty\) after taking the \(N\) and \(\epsilon\) limits.
This proves the local upper bound.
A finite cover of the compact simplex by the resulting
neighbourhoods gives the upper bound for every closed \(F\).

For the lower bound first suppose all \(x_r>0\).
Tilt the law of cycles by \(\prod_rx_r^{n_r(C)}\), and denote
the resulting law by \(\nu_{N,x}\).
Its limiting logarithmic moment-generating function is
\[
 \log\sum_rx_re^{\theta_r},
\]
by applying the partition-function formula at the two fixed
positive weight vectors \(x\) and \((x_re^{\theta_r})_r\).
For each coordinate, a sufficiently small fixed positive or
negative tilt and Chernoff's inequality show
\[
 \nu_{N,x}(B_\epsilon(x))=1-\exp(-\Omega_{x,\epsilon}(N)).
\]
For a constant coordinate the corresponding deviation event
is empty. There are only finitely many coordinates.
On \(B_\epsilon(x)\),
\[
 \left|\sum_r(X_{N,r}-x_r)\log x_r\right|
                        \leq\epsilon\sum_r|\log x_r|.
\]
The exact Radon--Nikodym identity is
\[
 \mathbb P(X_N\in B_\epsilon(x))
 =\frac{W_N(x)}{W_N(\one)}
   \mathbb E_{\nu_{N,x}}\left[
       \exp\left(-N\sum_r X_{N,r}\log x_r\right)
                        \one_{B_\epsilon(x)}(X_N)\right].
\]
The pressure ratio has logarithm \(-N\log k+o_x(N)\).
Consequently the lower limiting exponent is at least
\(-I_R(x)-\epsilon\sum_r|\log x_r|\).
Send \(\epsilon\) to zero to obtain the local lower bound.

Every boundary point of the simplex is a limit of interior
points, and \(I_R\) is continuous on the whole simplex.
Apply the interior lower bound in a smaller ball about an
interior point contained in the prescribed neighbourhood.
This proves the boundary local lower bound as well.
For an open \(O\), choose any such ball contained in \(O\),
and optimize its centre. This proves the stated principle,
without requiring a theorem uniform in the tilt vector.
\end{proof}

\begin{corollary}[Exponential concentration and profile counts]
\label{cor:primeprofilecount}
For every fixed \(\epsilon>0\) there is \(c_{q,R,\epsilon}>0\)
such that a uniform prime circle in \(G_N(q,R)\) satisfies
\[
 \mathbb P\left\{\max_{r\in R}
          \left|\frac{n_r(C_N)}N-\frac1k\right|\geq\epsilon\right\}
           \leq\exp(-c_{q,R,\epsilon}N)
\]
for all sufficiently large even \(N\), unless the event is empty.
For a prescribed approximate profile \(x\), its fraction of
all prime circles has the exact exponential cost \(I_R(x)\)
in the iterated neighbourhood-limit sense of the theorem.
Thus the total count and the empirical-profile count are
determined at exponential precision.
\end{corollary}
\begin{proof}
The rate is nonnegative and vanishes only at the uniform
vector. Its infimum on the indicated closed set is strictly
positive, whenever that set is nonempty.
The large-deviation upper bound gives the displayed estimate
with any fixed smaller positive exponent.
The local formula gives the remaining assertion.
\end{proof}

\begin{remark}[The edges are constrained, not independent]
Each prime circle obeys the exact identity
\[
 \sum_{r\in R}r\,n_r(C)
       \equiv\sum_{ij\in E(C)}(i+j)
       =2\sum_{i=1}^N i=N(N+1)\pmod q.
\]
The large-deviation law agrees with independent uniform
residue sampling only at the exponential profile scale.
It does not assert independence, a central limit theorem,
or positive counts for every exact integer profile.
For example \(q=8,R=\{1,3\}\) satisfies \(g_R=2\).
The number \(n_3(C)\) always obeys
\[
 n_3(C)\equiv N^2/2\pmod4.
\]
Nevertheless its normalized proportion has binary rate
\[
 I(x)=x\log(2x)+(1-x)\log(2(1-x)).
\]
An exact profile using only one class has no Hamilton cycle
for large \(N\), by the obstruction theorem.
A fixed positive-width neighbourhood of that boundary
profile has the finite exponential rate \(I(0)=I(1)=\log2\).
The order of limits is essential here.
\end{remark}
